\section{Method}
\label{sec:method}

We present a budgeted corpus selector that assembles compact evidence-bearing subgraphs from naturally linked document collections. The method operates in three stages: dense seed retrieval, link-context graph walking, and budget-aware fill. All stages are governed by an explicit selector budget that bounds the retained corpus mass of the selected subgraph. Figure~\ref{fig:architecture} illustrates the overall pipeline.

\subsection{Overview}
\label{sec:method:overview}

Given a query $q$ and a link-context graph $\mathcal{G} = (\mathcal{V}, \mathcal{E})$ as defined in Section~\ref{sec:problem}, our corpus selector $\sigma$ produces a selected set $S \subseteq \mathcal{V}$ through the following stages:

\begin{enumerate}[leftmargin=*,nosep]
\item \textbf{Seed retrieval.} A dense retriever identifies up to $k$ starting documents from $\mathcal{V}$ by embedding similarity to $q$, forming the seed set $S_0$.

\item \textbf{Link-context walk.} Starting from $S_0$, the selector traverses edges in $\mathcal{G}$ for up to $H$ expansion steps. At each step, outgoing hyperlinks from the current document are scored using their link context $\ell_e = (a_e, c_e, i_e)$, and the highest-scoring neighbor is added to the selected set. The walk terminates when the step budget $H$ is exhausted, no unvisited neighbors remain, or (in the controller variant) the scorer signals a stop.

\item \textbf{Budget fill.} After the walk terminates, any remaining selector headroom is filled by adding the highest-ranked unselected documents from a dense retrieval pool, subject to a quality gate. This ensures high budget utilization without inflating the walk itself.
\end{enumerate}

The complete selector budget is the triple $(\rho, H, k)$ in the canonical ratio-controlled setting. The method is zero-shot: no component requires task-specific training or fine-tuning.

\subsection{Graph Representation}
\label{sec:method:graph}

The input graph $\mathcal{G} = (\mathcal{V}, \mathcal{E})$ is a \emph{link-context graph}. Each node $v \in \mathcal{V}$ represents a document with title $t_v$ and sentence sequence $\mathbf{s}_v$, as formalized in Section~\ref{sec:problem:graph}. Each directed edge $e = (u, v) \in \mathcal{E}$ corresponds to a hyperlink from document $u$ to document $v$ and carries a link-context triple $\ell_e = (a_e, c_e, i_e)$: anchor text $a_e$, the sentence $c_e$ in $u$ containing the anchor, and its sentence index $i_e$.

This representation differs from knowledge graphs in two ways. First, it requires no entity extraction or relation typing---edges are the hyperlinks already present in the corpus. Second, it preserves full document structure rather than decomposing text into entity-centric triples. The link context provides a local semantic signal about each neighbor's relevance \emph{before} the neighbor's content is retrieved, enabling informed traversal decisions without reading the target document.

Multiple edges may exist between the same document pair (corresponding to distinct hyperlinks), each with its own context triple. The graph is constructed directly from corpus metadata without any corpus-wide preprocessing beyond link extraction.

\subsection{Stage 1: Dense Seed Retrieval}
\label{sec:method:seed}

The walk requires one or more starting documents. We obtain these through dense retrieval using a sentence-transformer model~\cite{reimers2019sentencebert}. Specifically, we encode the query $q$ and each document $v \in \mathcal{V}$ (using the concatenation of title and text) with \stmodel{}, then rank documents by cosine similarity to $q$ and select the top-$k$ as the seed set $S_0$.

Dense seeding serves two roles. As a standalone baseline (\seldense{}), top-$k$ retrieval under the same selector budget provides the flat dense control. As the first stage of any walk selector, it anchors the walk in the most query-relevant region of the graph.

\subsection{Stage 2: Link-Context Walk}
\label{sec:method:walk}

From the seed set $S_0$, the selector traverses hyperlinks in $\mathcal{G}$ to discover additional evidence. We implement walk strategies of increasing sophistication, all operating over the link-context graph from a shared seed set.

\paragraph{Single-path walk.}
The simplest strategy maintains a single walk frontier. Starting from the highest-scoring seed in $S_0$, the walker iterates: at each step $t \in \{1, \ldots, H\}$, it collects all outgoing edges from the current document to unvisited neighbors, scores each edge using its link context (Section~\ref{sec:method:scoring}), and moves to the neighbor with the highest score. The walk terminates when the step budget $H$ is exhausted, no unvisited neighbors exist (\deadend{}), or the best score falls below a minimum threshold (\scorebelow{}). Algorithm~\ref{alg:single-path} summarizes this procedure.

\begin{algorithm}[t]
\caption{Single-path link-context walk}
\label{alg:single-path}
\begin{algorithmic}[1]
\Require Query $q$, graph $\mathcal{G}$, seed set $S_0$, step budget $H$, scorer $f$, threshold $\theta$
\State $v_0 \gets \arg\max_{v \in S_0} f_{\text{node}}(q, v)$; \quad $S \gets \{v_0\}$; \quad $t \gets 0$
\While{$t < H$}
  \State $C \gets \{(u, v, \ell) \in \mathcal{E} : u = v_t,\; v \notin S\}$
  \If{$C = \emptyset$}
    \State \Return $S$ \Comment{dead end}
  \EndIf
  \State $(u^*, v^*, \ell^*) \gets \arg\max_{(u,v,\ell) \in C} f(q, \ell, \mathcal{G}, S)$
  \If{$f(q, \ell^*, \mathcal{G}, S) < \theta$}
    \State \Return $S$ \Comment{score below threshold}
  \EndIf
  \State $v_{t+1} \gets v^*$; \quad $S \gets S \cup \{v^*\}$; \quad $t \gets t + 1$
\EndWhile
\State \Return $S$
\end{algorithmic}
\end{algorithm}

\paragraph{Constrained multi-path walk.}
This variant extends the single-path walk with an LLM-guided controller that can take three actions at each step: \actchooseone{} (continue on the primary link), \actchoosetwo{} (fork a scout branch to a secondary link), or \actstop{} (terminate the walk). The controller emits per-candidate utility scores decomposed into \emph{direct support}, \emph{bridge potential}, and \emph{redundancy risk}, along with an \emph{evidence cluster confidence} that governs early stopping and budget pacing.

Branching is explicitly bounded: at most one fork is allowed per walk, and the scout branch contributes its target document to the selected set only if its token cost falls within 10\% of the total budget and its redundancy risk is below 0.60. This constraint prevents the precision collapse observed in unconstrained branchy search. A single backtrack is also permitted: if the current path reaches a dead end or low-scoring region, the walker reverts one step and follows the controller's designated backup edge.

\paragraph{Baseline walk strategies.}
For comparison, we implement two dense-only iterative baselines that do not use hyperlink structure:

\begin{itemize}[leftmargin=*,nosep]
\item \seliterdense{}: at each hop, the query is augmented with the title and lead sentence of the current frontier documents, and a new dense retrieval is performed over the remaining corpus. This follows the retrieve-then-expand pattern of MDR~\cite{xiong2021mdr}.
\item \selmdrlight{}: each frontier document independently issues its own dense retrieval query, and the per-frontier results are merged by taking the highest score per candidate. This is a repo-native approximation of per-hop MDR without trained query encoders.
\end{itemize}

Both baselines operate under the same selector budget and seed retrieval as the walk selectors, providing a controlled comparison surface.

\subsection{Edge Scoring}
\label{sec:method:scoring}

At each walk step, outgoing edges are scored to decide which neighbor to visit. We implement three scorer families, all operating on the link-context triple $\ell_e = (a_e, c_e, i_e)$ \emph{without reading the target document's content}. This is a key property: the scorer makes navigation decisions based on local information available at the source document, avoiding the cost of pre-fetching candidate pages.

\paragraph{Lexical overlap scoring.}
The overlap scorer computes a weighted combination of normalized token overlap between the query $q$ and each component of the link context:
\begin{equation}
\label{eq:overlap}
f_{\text{ov}}(q, \ell_e) = w_a \cdot \phi(q, a_e) + w_c \cdot \phi(q, c_e) + w_t \cdot \phi(q, t_{v}) + w_n \cdot \mathbb{1}[v \notin S],
\end{equation}
where $\phi(\cdot, \cdot)$ is the Jaccard-style normalized token overlap after stopword removal, $t_v$ is the target document's title, and the novelty term $w_n$ rewards unvisited targets. Named weight profiles control the tradeoff between anchor and context signals (e.g., \texttt{overlap\_balanced}: $w_a{=}0.60, w_c{=}0.40$; \texttt{overlap\_title\_aware}: $w_a{=}0.45, w_c{=}0.35, w_t{=}0.20$).

\paragraph{Embedding-based scoring.}
The sentence-transformer scorer encodes the query and a textual representation of each edge (combining anchor text, sentence context, and target title) using the same \stmodel{} model as seed retrieval. The edge score is:
\begin{equation}
\label{eq:st}
f_{\text{ST}}(q, \ell_e) = w_d \cdot \cos(\mathbf{q}, \mathbf{e}) + w_n \cdot \mathbb{1}[v \notin S].
\end{equation}
When lookahead is enabled ($L > 1$), the scorer also computes the best one-hop-ahead score by encoding all edges reachable from the target $v$, without retrieving any document content:
\begin{equation}
\label{eq:st-lookahead}
f_{\text{ST+LA}}(q, \ell_e) = w_d \cdot \cos(\mathbf{q}, \mathbf{e}) + w_f \cdot \max_{e' \in \text{out}(v)} \cos(\mathbf{q}, \mathbf{e}') + w_n \cdot \mathbb{1}[v \notin S].
\end{equation}
Named profiles govern the balance of direct evidence versus future traversal promise (e.g., \texttt{st\_balanced}: $w_d{=}0.55, w_f{=}0.35, w_n{=}0.10$; \texttt{st\_future\_heavy}: $w_d{=}0.45, w_f{=}0.45, w_n{=}0.10$). The best non-LLM configuration in our experiments uses the \texttt{st\_future\_heavy} profile with lookahead depth $L{=}2$.

\paragraph{LLM-based scoring.}
The LLM scorer presents the query, the current path context, and a pre-filtered set of candidate edges to a language model, which returns per-edge relevance judgments. In the controller variant (used by \selcontroller{}), the LLM additionally selects an action (\actstop{}, \actchooseone{}, \actchoosetwo{}) and decomposes utility into subscores for direct support, bridge potential, and redundancy risk. To manage the candidate set presented to the LLM, a two-stage prefilter narrows outgoing edges: first by lexical overlap, then by semantic similarity, with a rescue mechanism that preserves candidates whose anchor text closely matches query entities. The two-stage prefilter uses lightweight embedding and lexical scorers. In the controller variant, the LLM receives a narrowed candidate set; we use \controllerlm{} as the primary controller (Section~\ref{sec:setup}) and evaluate sensitivity to model capability in Section~\ref{sec:analysis}.

\subsection{Stage 3: Budget Enforcement and Fill}
\label{sec:method:budget}

Budget enforcement operates at two levels.

\paragraph{Walk-level budget.}
The walk budget constrains the expansion by a maximum number of steps $H$, a minimum score threshold $\theta$, and a revisit policy. Each step adds exactly one document to the selected set, so the walk produces at most $H + |S_0|$ documents. The walk does not directly track token counts; token-level enforcement is deferred to the selection-level budget.

\paragraph{Selection-level budget.}
The selector budget enforces the hard cost constraint $\sum_{v \in S} \tau(v) \leq \rho T$. After the walk terminates with a candidate set $S_{\text{walk}}$, the selector trims documents in reverse selection order until the constraint is satisfied. For iterative-dense baselines, documents are added greedily in selection order until the budget is exhausted.

\paragraph{Budget fill with relative drop.}
When the walk terminates with remaining headroom, a post-walk fill stage retrieves the top-$K$ unselected documents (default $K{=}64$) by dense similarity from a backfill pool and adds them greedily until the selector budget is saturated. To prevent adding low-quality filler, we apply the \emph{relative drop} quality gate: candidates from the backfill pool are ranked by dense similarity to the query, and each candidate $v_i$ (in rank order $i = 1, 2, \ldots$) is admitted only while its similarity remains above a fraction $\rho_{\text{fill}}$ (default $0.5$) of the \emph{top-ranked} backfill candidate's score:
\begin{equation}
\label{eq:budget-fill}
\text{admit}(v_i) = \mathbb{1}\!\left[\text{sim}(q, v_i) \geq \rho_{\text{fill}} \cdot \text{sim}(q, v_1)\right],
\end{equation}
where $v_1$ is the highest-scoring unselected candidate in the backfill pool (not a walk step index) and $v_i$ is the $i$-th ranked candidate under consideration. The gate therefore compares each candidate to the best available backfill document, not to itself. This mechanism eliminates empty selections entirely---achieving a $0.0$ empty rate across all tested configurations---and raises budget utilization without diluting precision, since only documents sufficiently similar to the query are admitted.

Budget utilization, defined as $\sum_{v \in S} \tau(v) / (\rho T)$, is reported alongside selection quality. High utilization under high support F1 indicates efficient use of the selector budget; low utilization suggests premature termination. Budget adherence $\mathbb{1}[\sum_{v \in S} \tau(v) \leq \rho T]$ is verified as a hard constraint for every selector configuration.

\subsection{Selector Family}
\label{sec:method:family}

Table~\ref{tab:selector-family} summarizes the selector configurations compared in our experiments. Not all are claimed as contributions; the family is designed as a comparison class with a shared budget accounting framework.

\begin{table}[t]
\centering
\small
\caption{Selector family overview. All selectors share the same seed retrieval and budget accounting. The \emph{Walk} column indicates whether the selector traverses hyperlinks.}
\label{tab:selector-family}
\begin{tabular}{@{}lccl@{}}
\toprule
\textbf{Selector} & \textbf{Walk} & \textbf{LLM} & \textbf{Role} \\
\midrule
\seldense & \ding{55} & \ding{55} & Flat dense control \\
\seliterdense & \ding{55} & \ding{55} & Iterative dense baseline \\
\selmdrlight & \ding{55} & \ding{55} & Per-hop dense baseline \\
\selwalk & \ding{51} & \ding{55} & Hyperlink-local walk \\
\selcontroller & \ding{51} & \ding{51} & Controller-guided walk \\
\bottomrule
\end{tabular}
\end{table}

The \seldense{} selector serves as the primary control: it applies top-$k$ sentence-transformer retrieval under the same budget, with budget fill, but no graph traversal. The \seliterdense{} and \selmdrlight{} selectors are repo-native iterative-dense baselines that follow the same retrieve-and-expand pattern as MDR~\cite{xiong2021mdr} but without trained query encoders. The \selwalk{} selector is the main non-LLM method: a greedy hyperlink-local walk with link-context scoring. The \selcontroller{} selector augments the walk with an LLM controller and represents the main method candidate. All selectors compose with \budgetfill{} to ensure comparable budget utilization.
